\documentclass[10pt,a4paper]{amsart}
\usepackage[margin=19mm]{geometry}
\usepackage{amsmath,amssymb,mathtools}
\usepackage[scr=rsfs,cal=euler]{mathalfa}
\usepackage{xfrac}
\usepackage[unicode,hidelinks,bookmarks=false]{hyperref}
\newcommand{\ie}{i.e.\ }
\newcommand{\cf}{cf.\ }
\newcommand{\resp}{respectively\ }
\newcommand{\Subsection}{\subsection}
\providecommand{\nobreakdash}{\nobreak\hbox{-}}
\setlength{\emergencystretch}{2em}
\multlinegap0pt
\usepackage{tikz,booktabs,bm}
\usepackage{pgfplots}
\usepackage{tikz-3dplot}
\usepackage{times}
\usetikzlibrary{positioning, arrows.meta, decorations.pathmorphing, calc}
\pgfplotsset{compat=1.18}
\def\b{\beta}
\def\g{\gamma}
\def\Dn{\Delta^{(n)}}
\def\Dtherm{\mathrm{D}_{\mathrm{th}}}
\newcommand{\EMD}{\mathrm{EMD}}
\newcommand{\HMC}{\mathrm{HMS}}
\newcommand{\NMR}{\mathrm{NMR}}
\newcommand{\W}{W_1}
\newcommand{\Rn}{R^{(n)}_t}
\newcommand{\Rinf}{R^{(\infty)}_t}
\newcommand{\Tn}{T^{(n)}_t}
\newcommand{\Tinf}{T^{(\infty)}_t}
\newcommand{\be}{\beta}
\newcommand{\bc}{\beta_c}
\newcommand{\sr}{\varrho}      % spectral radius
\newtheorem{theorem}{Theorem}[section]
\newtheorem{proposition}[theorem]{Proposition}
\newtheorem{lemma}[theorem]{Lemma}
\newtheorem{corollary}[theorem]{Corollary}
\newtheorem{definition}[theorem]{Definition}
\newtheorem{remark}[theorem]{Remark}
\begin{document}
\title[Weight geometry governs functional memory in complex systems]{Weight geometry governs functional memory in complex systems}
\author{Elkaïoum M. Moutuou, Habib Benali}
\date{}
\maketitle
\noindent\emph{Source and licence.} Weight geometry governs functional memory in complex systems. Version arXiv:2606.25826v2 (CC BY 4.0). This material is distributed under the Creative Commons Attribution 4.0 International licence, http://creativecommons.org/licenses/by/4.0/, as recorded in the arXiv OAI licence metadata for this exact version and shown on the abstract page https://arxiv.org/abs/2606.25826v2. Full rights notice: licenses/arxiv-2606.25826-CC-BY-4.0.txt.\par\smallskip
\noindent\textit{Editorial scope.} Counted unit: the original \texttt{lemma} environment \texttt{lem:monotone} at source lines 1711--1727 together with its complete proof at lines 1729--1754; the delivered document keeps source lines 1697--1755 so that every referenced label and every citation resolves inside it. Native editable source is the paper's own concatenated TeX; the AMS wrapper is editorial formatting only. No publisher class, upstream script or unrelated artwork is redistributed.
\section{Monotonicity and normalization of the thermalization cascade}
\label{sec:monotonicity}
% =====================================================================

Let $A$ be a non-negative adjacency matrix with spectral radius $\rho(A)$.
We work in the regime $\beta>\beta_c=\log\rho(A)$, so that
$e^{-\beta}\rho(A)<1$ and the Neumann series
\[
  R_\beta^{(\infty)}
  =
  \sum_{k=0}^{\infty}e^{-\beta k}A^k
\]
converges entrywise and in any matrix norm.


\begin{lemma}[Monotonicity in memory depth]
\label{lem:monotone}
For every $\beta>\beta_c$,
\[
  R_\beta^{(\infty)}-R_\beta^{(n+1)}
  \;\leq\;
  R_\beta^{(\infty)}-R_\beta^{(n)}
  \qquad\text{entrywise,}
\]
and consequently $E_{n+1}(\beta)\leq E_n(\beta)$, where
\[
  E_n(\beta)
  =
  \frac{\|R_\beta^{(\infty)}-R_\beta^{(n)}\|_F}
       {\|R_\beta^{(\infty)}\|_F}.
\]
\end{lemma}

\begin{proof}
Since $A$ is non-negative, the tail sum
\[
  R_\beta^{(\infty)}-R_\beta^{(n)}
  =
  \sum_{k=n+1}^{\infty}e^{-\beta k}A^k
\]
is an entrywise non-negative matrix. Writing
\[
  R_\beta^{(\infty)}-R_\beta^{(n)}
  =
  e^{-\beta(n+1)}A^{n+1}
  +
  \bigl(R_\beta^{(\infty)}-R_\beta^{(n+1)}\bigr),
\]
and noting that $e^{-\beta(n+1)}A^{n+1}\geq 0$ entrywise, we conclude
\[
  R_\beta^{(\infty)}-R_\beta^{(n+1)}
  \;\leq\;
  R_\beta^{(\infty)}-R_\beta^{(n)}
  \qquad\text{entrywise.}
\]
Since the Frobenius norm is monotone on entrywise non-negative matrices,
dividing by $\|R_\beta^{(\infty)}\|_F>0$ gives $E_{n+1}(\beta)\leq
E_n(\beta)$.
\end{proof}


\end{document}
